%% LyX 2.4.4 created this file.  For more info, see https://www.lyx.org/.
%% Do not edit unless you really know what you are doing.
\documentclass[english]{beamer}
\usepackage[T1]{fontenc}
\usepackage[latin9]{inputenc}
\setcounter{secnumdepth}{3}
\setcounter{tocdepth}{3}
\PassOptionsToPackage{normalem}{ulem}
\usepackage{ulem}

\makeatletter
%%%%%%%%%%%%%%%%%%%%%%%%%%%%%% Textclass specific LaTeX commands.
% this default might be overridden by plain title style
\newcommand\makebeamertitle{\frame{\maketitle}}%
% (ERT) argument for the TOC
\AtBeginDocument{%
  \let\origtableofcontents=\tableofcontents
  \def\tableofcontents{\@ifnextchar[{\origtableofcontents}{\gobbletableofcontents}}
  \def\gobbletableofcontents#1{\origtableofcontents}
}

%%%%%%%%%%%%%%%%%%%%%%%%%%%%%% User specified LaTeX commands.
\beamertemplateshadingbackground{red!5}{structure!5}

\usepackage{beamerthemeshadow}
\usepackage{pgfnodes,pgfarrows,pgfheaps}
\usepackage{caption}

\beamertemplatetransparentcovereddynamicmedium

\pgfdeclareimage[width=0.6cm]{icsi-logo}{beamer-icsi-logo}
\logo{\pgfuseimage{icsi-logo}}

\newcommand{\Class}[1]{\operatorname{\mathchoice
  {\text{\small #1}}
  {\text{\small #1}}
  {\text{#1}}
  {\text{#1}}}}

\newcommand{\Lang}[1]{\operatorname{\text{\textsc{#1}}}}

% This gets defined by beamerbasecolor.sty, but only at the beginning of
% the document
\colorlet{averagebackgroundcolor}{normal text.bg}

\newcommand{\tape}[3]{%
  \color{structure!30!averagebackgroundcolor}
  \pgfmoveto{\pgfxy(-0.5,0)}
  \pgflineto{\pgfxy(-0.6,0.1)}
  \pgflineto{\pgfxy(-0.4,0.2)}
  \pgflineto{\pgfxy(-0.6,0.3)}
  \pgflineto{\pgfxy(-0.4,0.4)}
  \pgflineto{\pgfxy(-0.5,0.5)}
  \pgflineto{\pgfxy(4,0.5)}
  \pgflineto{\pgfxy(4.1,0.4)}
  \pgflineto{\pgfxy(3.9,0.3)}
  \pgflineto{\pgfxy(4.1,0.2)}
  \pgflineto{\pgfxy(3.9,0.1)}
  \pgflineto{\pgfxy(4,0)}
  \pgfclosepath
  \pgffill

  \color{structure}  
  \pgfputat{\pgfxy(0,0.7)}{\pgfbox[left,base]{#1}}
  \pgfputat{\pgfxy(0,-0.1)}{\pgfbox[left,top]{#2}}

  \color{black}
  \pgfputat{\pgfxy(-.1,0.25)}{\pgfbox[left,center]{\texttt{#3}}}%
}

\newcommand{\shorttape}[3]{%
  \color{structure!30!averagebackgroundcolor}
  \pgfmoveto{\pgfxy(-0.5,0)}
  \pgflineto{\pgfxy(-0.6,0.1)}
  \pgflineto{\pgfxy(-0.4,0.2)}
  \pgflineto{\pgfxy(-0.6,0.3)}
  \pgflineto{\pgfxy(-0.4,0.4)}
  \pgflineto{\pgfxy(-0.5,0.5)}
  \pgflineto{\pgfxy(1,0.5)}
  \pgflineto{\pgfxy(1.1,0.4)}
  \pgflineto{\pgfxy(0.9,0.3)}
  \pgflineto{\pgfxy(1.1,0.2)}
  \pgflineto{\pgfxy(0.9,0.1)}
  \pgflineto{\pgfxy(1,0)}
  \pgfclosepath
  \pgffill

  \color{structure}  
  \pgfputat{\pgfxy(0.25,0.7)}{\pgfbox[center,base]{#1}}
  \pgfputat{\pgfxy(0.25,-0.1)}{\pgfbox[center,top]{#2}}

  \color{black}
  \pgfputat{\pgfxy(-.1,0.25)}{\pgfbox[left,center]{\texttt{#3}}}%
}

\pgfdeclareverticalshading{heap1}{\the\paperwidth}%
  {color(0pt)=(black); color(1cm)=(structure!65!white)}
\pgfdeclareverticalshading{heap2}{\the\paperwidth}%
  {color(0pt)=(black); color(1cm)=(structure!55!white)}
\pgfdeclareverticalshading{heap3}{\the\paperwidth}%
  {color(0pt)=(black); color(1cm)=(structure!45!white)}
\pgfdeclareverticalshading{heap4}{\the\paperwidth}%
  {color(0pt)=(black); color(1cm)=(structure!35!white)}
\pgfdeclareverticalshading{heap5}{\the\paperwidth}%
  {color(0pt)=(black); color(1cm)=(structure!25!white)}
\pgfdeclareverticalshading{heap6}{\the\paperwidth}%
  {color(0pt)=(black); color(1cm)=(red!35!white)}

\newcommand{\heap}[5]{%
  \begin{pgfscope}
    \color{#4}
    \pgfheappath{\pgfxy(0,#1)}{\pgfxy(-#2,0)}{\pgfxy(#2,0)}
    \pgfclip
    \begin{pgfmagnify}{1}{#1}
      \pgfputat{\pgfpoint{-.5\paperwidth}{0pt}}{\pgfbox[left,base]{\pgfuseshading{heap#5}}}
    \end{pgfmagnify}
  \end{pgfscope}
  %\pgffill
  
  \color{#4}
  \pgfheappath{\pgfxy(0,#1)}{\pgfxy(-#2,0)}{\pgfxy(#2,0)}
  \pgfstroke

  \color{white}
  \pgfheaplabel{\pgfxy(0,#1)}{#3}%
}


\pgfdeclareradialshading{graphnode}
  {\pgfpoint{-3pt}{3.6pt}}%
  {color(0cm)=(beamerexample!15);
    color(2.63pt)=(beamerexample!75);
    color(5.26pt)=(beamerexample!70!black);
    color(7.6pt)=(beamerexample!50!black);
    color(8pt)=(beamerexample!10!averagebackgroundcolor)}

\newcommand{\graphnode}[2]{
  \pgfnodecircle{#1}[virtual]{#2}{8pt}
  \pgfputat{#2}{\pgfbox[center,center]{\pgfuseshading{graphnode}}}
}

%\setbeamerfont{author}{size=\Huge} 
%\setbeamerfont{date}{size=\small} 
%\setbeamerfont{exampletext}{series=\bfseries}
%\setbeamerfont{frametitle}{series=\bfseries,size=\large}
%\setbeamerfont{framesubtitle}{series=\mdseries,size=\large}
%\setbeamerfont{section}{series=\bfseries} 
\setbeamerfont{title}{series=\bfseries}

\setbeamercovered{invisible}

\makeatother

\usepackage{babel}
\begin{document}
\title{College Algebra}
\subtitle{MTH 131~~|~~Section B and I~~|~~Fall 2012}
\institute{{\Large\textbf{Johnson C. Smith }}}
\author{Dr.\,James Ashe}
\date{{}}

\makebeamertitle
%%suppress auto date
\title{Algebra: Review \& Preview}

\subtitle{}

\author{}

\institute{}

\makebeamertitle
\begin{frame}{Order of operations}

\begin{columns}


\column{6cm}
\begin{itemize}
\item <2->\textbf{\uline{P}}lease: Parenthesis and brackets
\item <3->\textbf{\uline{E}}xcuse: Exponents and Radicals
\item <4->\textbf{\uline{M}}y: Multiplication
\item <5->\textbf{\uline{D}}ear: Division
\item <6->\textbf{\uline{A}}unt: Addition
\item <7->\textbf{\uline{S}}ally: Subtraction
\end{itemize}

\column{6cm}
\begin{itemize}
\item []<8->$2(4-6)^{2}+\frac{8+1}{3}-1$
\item []<9->$=2(4-6)^{2}+\frac{(8+1)}{3}-1$
\item []<10->$=2(-2)^{2}+\frac{9}{3}-1$
\item []<11->$=2(4)+\frac{9}{3}-1$
\item []<12->$=8+\frac{9}{3}-1$
\item []<13->$=8+3-1$
\item []<14->$=11-1$
\item []<15->$=10$
\end{itemize}
\end{columns}
\end{frame}

\begin{frame}{}

Note that subtraction can be rewritten as addition.
\begin{itemize}
\item <2->[ex)]$4-3=$\uncover<3->{$4+(-3)=1$}
\end{itemize}
\uncover<4->{Also, division can be rewritten as multiplication.}
\begin{itemize}
\item <5->[ex)]$\frac{9}{3}=$\uncover<6->{$\frac{1}{3}\cdot9=3$}
\end{itemize}
\uncover<7->{So the acronym could've been P.E.D.M.S.A.}
\end{frame}

\begin{frame}{FRACTIONS!!!!}

\begin{itemize}
\item [ex)]${\color{blue}{\color{blue}}\frac{1}{2}+}\mathinner{\color{blue}\frac{3}{4}}\mathrel{\color{red}\neq}\frac{{\color{blue}1+3}}{{\color{blue}2+4}}$\uncover<2->{\textrm{\textcolor{blue}{$=\frac{4}{6}=\frac{2}{3}$}}\textrm{}}
\item <3->[]But why?
\item <4->[]\textcolor{blue}{$\frac{1}{2}+\frac{3}{4}=\frac{1}{2}(1)+\frac{1}{4}(3)$}\uncover<5->{
\\
You have to do the division/multiplication steps BEFORE the addition.}
\item <6->[ex)]\textcolor{blue}{$\frac{1}{2}(1)+\frac{3}{4}$}\uncover<7->{$\mathrel{\color{blue}=}{\color{blue}\frac{1}{2}(\frac{2}{2})+}\mathinner{\color{blue}\frac{3}{4}}$}\uncover<8->{$\mathrel{\color{blue}=}{\color{blue}\frac{2}{4}+}\mathinner{\color{blue}\frac{3}{4}}$}
\item <9->[]Now, \textcolor{blue}{$\frac{2}{4}+\frac{3}{4}=\frac{2+3}{4}=\frac{5}{4}$}
\item <10->[]But why are we allowed to do this now?
\item <11->[]\textcolor{blue}{$\frac{2}{4}+\frac{3}{4}$}\uncover<12->{\textcolor{blue}{$=\frac{1}{4}(2)+\frac{1}{4}(3)$}}\uncover<13->{$\mathrel{\color{blue}=}{\color{blue}\frac{1}{4}(2+3}\mathclose{\color{blue})}$}
\item <12->[]Now we \uline{have} to do the addition first!
\end{itemize}
\end{frame}

\begin{frame}{Solving Equations}

\end{frame}

\begin{frame}{Homework}

\textbf{HW pg. 86-87: 2, 24, 30, 32, 46, 90, 98, 110, 120 }
\end{frame}

\end{document}
